\documentclass[11pt,a4paper]{article}
\usepackage[T1]{fontenc}
\usepackage[utf8]{inputenc}
\usepackage{amsmath,amsthm,mathtools}
\usepackage{newpxtext,newpxmath}
\usepackage[scaled=.94]{sourcesanspro}
\usepackage[margin=25mm,headheight=15pt]{geometry}
\usepackage{microtype}
\usepackage{xcolor,booktabs,array,longtable,float}
\usepackage{enumitem}
\usepackage{listings}
\usepackage{fancyhdr,titlesec}
\usepackage[most]{tcolorbox}
\usepackage{hyperref}
\usepackage[nameinlink,noabbrev]{cleveref}
\definecolor{Ink}{HTML}{203449}
\definecolor{Accent}{HTML}{176B76}
\definecolor{Pale}{HTML}{EFF5F5}
\definecolor{Muted}{HTML}{526271}
\hypersetup{colorlinks=true,linkcolor=Accent,citecolor=Accent,urlcolor=Accent,
  pdftitle={Binary Runs in Rueppel Hankel Determinants},
  pdfauthor={Research report prepared for Vladimir Reshetnikov},
  pdfsubject={Four Rueppel Hankel conjectures, a polynomial refinement, and a correction}}
\titleformat{\section}{\Large\sffamily\bfseries\color{Ink}}{\thesection}{.65em}{}
\titleformat{\subsection}{\large\sffamily\bfseries\color{Accent}}{\thesubsection}{.6em}{}
\titlespacing*{\section}{0pt}{2.0ex plus 1ex minus .2ex}{1.0ex plus .2ex}
\titlespacing*{\subsection}{0pt}{1.6ex plus .6ex}{.6ex}
\pagestyle{fancy}
\fancyhf{}
\fancyhead[L]{\small\sffamily\color{Muted}Rueppel Hankel determinants}
\fancyhead[R]{\small\sffamily\color{Muted}Research manuscript}
\fancyfoot[C]{\small\sffamily\thepage}
\renewcommand{\headrulewidth}{.3pt}
\setlength{\parskip}{.35em}
\setlength{\parindent}{1.2em}
\setlength{\emergencystretch}{2em}
\allowdisplaybreaks[2]
\numberwithin{equation}{section}
\newtheorem{theorem}{Theorem}[section]
\newtheorem{lemma}[theorem]{Lemma}
\newtheorem{proposition}[theorem]{Proposition}
\newtheorem{corollary}[theorem]{Corollary}
\theoremstyle{definition}
\newtheorem{definition}[theorem]{Definition}
\newtheorem{example}[theorem]{Example}
\theoremstyle{remark}
\newtheorem{remark}[theorem]{Remark}
\newcommand{\Z}{\mathbb Z}
\newcommand{\C}{\mathbb C}
\newcommand{\Q}{\mathbb Q}
\newcommand{\R}{\mathbb R}
\newcommand{\Runs}{\mathcal R}
\newcommand{\Ones}{\rho}
\newcommand{\eps}{\varepsilon}
\newcommand{\pop}{\operatorname{popcount}}
\newcommand{\Gray}{\operatorname{Gray}}
\newcommand{\ceil}[1]{\left\lceil #1\right\rceil}
\newcommand{\floor}[1]{\left\lfloor #1\right\rfloor}
\newcommand{\one}{\boldsymbol 1}
\newcommand{\ind}[1]{\boldsymbol{[}#1\boldsymbol{]}}
\newcommand{\oeis}[1]{\href{https://oeis.org/#1}{\texttt{#1}}}
\newcommand{\arxiv}[1]{\href{https://arxiv.org/abs/#1}{arXiv:#1}}
\newtcolorbox{resultbox}{colback=Pale,colframe=Accent,boxrule=.5pt,
  arc=1mm,left=3mm,right=3mm,top=2mm,bottom=2mm,breakable}
\lstset{language=Python,basicstyle=\small\ttfamily,keywordstyle=\color{Accent}\bfseries,
  commentstyle=\color{Muted},stringstyle=\color{Ink},breaklines=true,
  frame=single,rulecolor=\color{Accent!50},backgroundcolor=\color{Pale},
  columns=fullflexible,keepspaces=true,showstringspaces=false,
  aboveskip=1em,belowskip=.8em}
\setlist{nosep,leftmargin=1.5em}

\begin{document}
\hypersetup{pageanchor=false}
\begin{titlepage}
\vspace*{12mm}
{\sffamily\color{Accent}\large INTEGER SEQUENCES \enspace / \enspace EXACT DETERMINANTS\par}
\vspace{9mm}
{\sffamily\bfseries\color{Ink}\fontsize{28}{33}\selectfont
Binary Runs in Rueppel\par Hankel Determinants\par}
\vspace{6mm}
{\sffamily\large\color{Muted}Four conjectures, a polynomial refinement,\par and a correction\par}
\vspace{12mm}
{\large Prepared for Vladimir Reshetnikov\par}
{\color{Muted}19 September 2026\par}
\vspace{12mm}
\begin{abstract}
Let $r(x)=\sum_{j\ge0}x^{2^j-1}$, and let $H_n(f)$ denote the Hankel
determinant of order $n$ of the coefficients of $f$. We determine, as a
polynomial in $t$, every determinant of
\[
 f_t(x)=t x+\frac{1}{r(x^2)}.
\]
If $\Runs(n)$ is the number of binary runs of $n$ and $\eps(n)$ is the lowest
bit of its binary reflected Gray code, the answer is
\[
 H_n(f_t)=(-1)^{1+\eps(n)+\lfloor\Runs(n)/2\rfloor}
 \bigl(\eps(n)+(\Runs(n)-\eps(n))t^2\bigr),\qquad n\ge1.
\]
In particular, $|H_n(f_1)|=\Runs(n)$. Together with reciprocal-minor and
parity-block identities, this proves the statements numbered 8, 10, 13,
and 23 in Paul Barry's 2021 note on generalized Rueppel sequences.
We also classify all exceptional complex parameters, compute exact
signed distributions and dyadic sums, and give a counterexample to
Conjecture 22 as printed, with a valid signed replacement.
The proofs are algebraic and apply to every index. An accompanying
standard-library Python implementation independently checks 1,192
Rueppel-related determinants and supplies reproducible data.
\end{abstract}
\vfill
\begin{resultbox}
\small
\textbf{Research status.} This is an unrefereed, independently developed
manuscript, not a claim of established publication priority. The selected
statements retain their conjecture labels in the checked 2021 source.
A focused search did not locate a later resolution, but does not exclude
one. Related foundational identities are already present in earlier
work; the source comparison and the limits of the novelty claim are
made explicit in Sections~\ref{sec:problem} and~\ref{sec:status}.
\end{resultbox}
\end{titlepage}
\pagenumbering{roman}
\hypersetup{pageanchor=true}
\tableofcontents
\newpage
\pagenumbering{arabic}

\section{The problem and the main result}\label{sec:problem}

A sequence supported only at the numbers $2^j-1$ seems an unlikely source
of complicated determinants. Nevertheless, taking a reciprocal of its
generating function produces rapidly growing signed coefficients, while
its Hankel determinants retain a remarkably small, digit-controlled
structure. The main problem studied here is to identify that structure
and to explain why it survives the reciprocal operation.

Write
\begin{equation}\label{eq:rdef}
 r_n=\ind{n+1\text{ is a power of }2},\qquad
 r(x)=\sum_{n\ge0}r_nx^n=1+x+x^3+x^7+x^{15}+\cdots.
\end{equation}
This is the Rueppel sequence, \oeis{A036987}~\cite{oeisrueppel}. The functional equation
\begin{equation}\label{eq:rfunctional}
 r(x)=1+xr(x^2)
\end{equation}
is immediate from its support. In Barry's note~\cite{barry2021},
Conjecture 13 asks whether the Hankel transform of $x+1/r(x^2)$ is a
signed version of the binary-run sequence \oeis{A005811}, shifted by one
place. We make the sign explicit and evaluate a whole polynomial family.

For a power series $f(x)=\sum_{j\ge0}f_jx^j$, our convention is
\begin{equation}\label{eq:hdef}
 H_n(f)=\det(f_{i+j})_{0\le i,j<n},\qquad H_0(f)=1.
\end{equation}
Thus $n$ is a matrix size. Barry's Hankel-transform index is one less:
his term with index $j$ is our $H_{j+1}$. This convention matters in every
comparison below.

For $n>0$, let $\Runs(n)$ count the maximal runs of equal digits in the
usual binary expansion of $n$, with no leading zeros. Set $\Runs(0)=0$.
Let
\begin{equation}\label{eq:epsilon}
 g(n)=n\mathbin{\mathrm{xor}}\floor{n/2},\qquad
 \eps(n)=g(n)\bmod2.
\end{equation}
Here $g(n)$ is the binary reflected Gray code of $n$. Equivalently,
$\eps(n)=1$ for $n\equiv1,2\pmod4$, and it is zero otherwise.

\begin{theorem}[Polynomial refinement of the binary-run conjecture]\label{thm:main}
For every $n\ge1$, the following identity holds in $\Z[t]$:
\begin{equation}\label{eq:main}
\boxed{\quad
 H_n\!\left(t x+\frac1{r(x^2)}\right)
 =(-1)^{1+\eps(n)+\lfloor\Runs(n)/2\rfloor}
 \left(\eps(n)+(\Runs(n)-\eps(n))t^2\right).
 \quad}
\end{equation}
The empty determinant is $H_0=1$, separately. Consequently, the displayed
identity remains valid after any specialization of $t$ in a commutative
ring.
\end{theorem}

The proof is given in Section~\ref{sec:mainproof}. It uses a rank-two
block calculation, not induction on large reciprocal coefficients.

\begin{corollary}[Conjecture 13]\label{cor:c13}
For $f_1(x)=x+1/r(x^2)$,
\begin{equation}\label{eq:c13}
 |H_n(f_1)|=\Runs(n)\quad(n\ge1).
\end{equation}
In the source's indexing, the absolute value of term $j$ is
$\oeis{A005811}(j+1)$.
\end{corollary}
\begin{proof}
At $t=1$, the parenthesis in \eqref{eq:main} is $\Runs(n)>0$.
\end{proof}

\subsection{Scope of the result}

The targeted statements are from a single source, rather than a
collection of conjectures synthesized after observing the answer.
The source's conjecture numbers and our conclusions are as follows.

\begin{center}
\small
\begin{tabular}{@{}c >{\raggedright\arraybackslash}p{.48\textwidth} >{\raggedright\arraybackslash}p{.30\textwidth}@{}}
\toprule
Number & Subject & Conclusion here\\
\midrule
8 & Binary-run parity in $H_{j+1}(1/(1+xr(x)))$ & Exact sign formula; proved.\\
10 & Hankel transform of $1-x/r(x^2)$ & Period $4$; proved.\\
13 & Hankel transform of $x+1/r(x^2)$ & Full polynomial refinement; proved.\\
23 & Product of unshifted and shifted Rueppel determinants & Exact binary-run parity; proved.\\
22 & Complement and first-difference determinants & False as printed; signed replacement proved.\\
\bottomrule
\end{tabular}
\end{center}

Allouche, Han, and Shallit~\cite{ahs} had already proved conjectures from
Barry's earlier 2020 paper~\cite{barry2020}, including results for
$1\pm xr(x)$ and recurrences for shifted Rueppel determinants. Those
results are part of the prior framework, not claimed as discoveries
here. We rederive the required ingredients so that the present argument
can be checked without importing their proofs. The four conjecture
numbers above refer specifically to the \emph{2021} paper, not to the
same numbers in the earlier paper.

\section{Binary statistics and formal-series conventions}\label{sec:digits}

All series are formal. No analytic convergence or choice of an inverse
function is involved in the expression $1/r(x^2)$: its constant term is
one, so its multiplicative reciprocal exists uniquely over $\Z$.

It is useful to distinguish two digit statistics. We write $\Runs(n)$
for the total number of constant-bit runs, and $\Ones(n)$ for the number
of runs consisting of ones. Thus, for example,
\[
 10=(1010)_2,\qquad \Runs(10)=4,\qquad \Ones(10)=2.
\]
We use the empty expansion at zero, so both statistics vanish there.
The sequences $\Runs(n)$ and $\Ones(n)\bmod2$ are respectively
\oeis{A005811} and \oeis{A268411}~\cite{oeisruns,oeisparity}.

\begin{lemma}[Digit identities]\label{lem:digits}
For $m\ge0$,
\begin{align}
 \Runs(2m)&=\Runs(m)+(m\bmod2),\label{eq:runs-even}\\
 \Runs(2m+1)&=\Runs(m)+1-(m\bmod2),\label{eq:runs-odd}\\
 \Ones(2m)&=\Ones(m),\label{eq:ones-even}\\
 \Ones(2m+1)&=\Ones(m)+1-(m\bmod2).\label{eq:ones-odd}
\end{align}
Moreover,
\begin{equation}\label{eq:grayruns}
 \Runs(n)=\pop(g(n)),\qquad
 \Ones(n)=\ceil{\Runs(n)/2}.
\end{equation}
If
\begin{equation}\label{eq:ab}
 a(n)=\pop\!\left(\floor{g(n)/2}\right),\qquad b(n)=g(n)\bmod2,
\end{equation}
then
\begin{equation}\label{eq:abidentities}
 b(n)=\eps(n),\quad a(n)=\Runs(n)-\eps(n)=\Runs(\floor{n/2}),
 \quad \binom n2\equiv a(n)\pmod2.
\end{equation}
\end{lemma}
\begin{proof}
Appending a zero creates a new run exactly when the previous last digit
is one. Appending a one creates a new run of ones, and a new total run,
exactly when the previous last digit is zero; the same statement gives
the stipulated value at $m=0$. This proves the first four identities.

Write the binary digits as $n_i$, with all sufficiently large $n_i=0$.
The $i$th Gray-code digit is $n_i\mathbin{\mathrm{xor}}n_{i+1}$.
A one in that code records a change between adjacent binary digits,
including the change from the leading one to the next higher zero.
Its number of ones therefore equals the number of binary runs. Since
the first run is a run of ones, runs of ones and zeros alternate, and
there are $\ceil{\Runs(n)/2}$ runs of ones.

Right shift commutes with the Gray-code map. Deleting the lowest
Gray-code digit thus leaves the Gray code of $\floor{n/2}$, proving the
first three assertions of \eqref{eq:abidentities}. Finally,
$\binom n2\bmod2$ is the second-lowest binary digit $n_1$. The xor of all
Gray-code digits with indices at least one telescopes to $n_1$.
That xor is the parity of their number of ones, namely $a(n)$.
\end{proof}

For shifted Hankel determinants we use
\begin{equation}\label{eq:shift}
 H_n^{[s]}(f)=\det(f_{i+j+s})_{0\le i,j<n},
 \qquad f_j=0\text{ when }j<0.
\end{equation}
Every determinant of size zero is one, whatever the shift. The
zero-extension convention in \eqref{eq:shift} will be needed for $s=-1$.

\section{Three determinant tools}\label{sec:tools}

\subsection{One coupling between two blocks}

If $A$ is a nonempty square matrix, let $A^{\circ}$ mean the matrix
obtained by deleting its first row and first column. In particular,
for a $1\times1$ matrix, $\det A^{\circ}=1$.

\begin{lemma}[Rank-two coupling]\label{lem:coupling}
Let $A$ and $B$ be square matrices of positive sizes, and let $e$ and $f$
be their first coordinate vectors. For any scalar $u$,
\begin{equation}\label{eq:coupling}
 \det\begin{pmatrix}A&u e f^T\\u f e^T&B\end{pmatrix}
 =\det A\det B-u^2\det A^{\circ}\det B^{\circ}.
\end{equation}
\end{lemma}
\begin{proof}
A term in the permutation expansion that uses neither off-diagonal
block belongs to $\det A\det B$. If a term uses one off-diagonal entry,
it must also use the other: the only available crossing entries join
the first coordinate of each block. These two entries form a
transposition, contributing $-u^2$; the remaining choices are exactly
the two deleted-block determinants. There are no other nonzero terms.
This proof requires no invertibility assumption.
\end{proof}

\subsection{Square and rectangular anti-diagonal blocks}

\begin{lemma}[Parity-block determinants]\label{lem:antiblock}
If $C$ is $m\times m$ and $A$ is any matrix of that size, then
\begin{equation}\label{eq:squareblock}
 \det\begin{pmatrix}A&C\\C^T&0\end{pmatrix}
 =(-1)^m(\det C)^2.
\end{equation}
If $C$ is $(m+1)\times m$, $e=(1,0,\ldots,0)^T$, and $C_*$ is obtained
by deleting the first row of $C$, then
\begin{equation}\label{eq:rectblock}
 \det\begin{pmatrix}u e e^T&C\\C^T&0\end{pmatrix}
 =u(-1)^m(\det C_*)^2.
\end{equation}
The second formula includes $m=0$.
\end{lemma}
\begin{proof}
In \eqref{eq:squareblock}, the last $m$ rows must use the first $m$
columns. Consequently the first $m$ rows use the last $m$ columns,
so $A$ contributes nothing. Exchanging the two column blocks gives
the sign $(-1)^{m^2}=(-1)^m$ and the product $\det C\det C^T$.

In \eqref{eq:rectblock}, there is one more upper row than lower column.
A nonzero permutation term must therefore use one entry from the
upper-left block. Its only possible entry is the first diagonal
entry $u$. Removing that row and column leaves the square-block
case with $C_*$. The corresponding cofactor has positive sign.
\end{proof}

\subsection{Complementary minors of reciprocal series}

The next identity is useful beyond the Rueppel sequence. Its negative
shift is the reason the apparently unrelated series $xr(x)$ enters the
main calculation.

\begin{lemma}[Reciprocal Hankel minors]\label{lem:reciprocal}
Let $A(x)=\sum_{j\ge0}a_jx^j$ have $a_0=1$, and put $C(x)=1/A(x)$.
For $n\ge0$ and $s\ge1$,
\begin{equation}\label{eq:recipgeneral}
 H_n^{[s]}(C)
 =(-1)^{n+\binom{s-1}{2}}\,
 H_{n+s-1}^{[2-s]}(A).
\end{equation}
Also, for $n\ge1$,
\begin{equation}\label{eq:recipordinary}
 H_n(C)=(-1)^{n-1}H_{n-1}^{[2]}(A).
\end{equation}
These are polynomial identities in the coefficients of $A$ and are
valid over every commutative ring.
\end{lemma}
\begin{proof}
For a series $A$, form the lower triangular Toeplitz matrix
$T_L(A)=(a_{i-j})_{0\le i,j<L}$. Its diagonal entries are one, and
formal series multiplication gives
\[
 T_L(C)=T_L(A)^{-1},\qquad \det T_L(A)=1.
\]
We first recall, with zero-based indices, the complementary-minor
identity
\begin{equation}\label{eq:jacobiminor}
 \det(M^{-1})_{I,J}
 =(-1)^{\sum I+\sum J}\frac{\det M_{J^c,I^c}}{\det M}.
\end{equation}
Here $I,J$ have the same size. To see it, permute the selected row and
column sets to the front and apply the block inverse formula. The
determinant of the relevant inverse block is the determinant of the
complementary block divided by the full determinant. The row and
column permutations together contribute $(-1)^{\sum I+\sum J}$.
A derivation assuming invertible blocks proves the identity for generic
matrices; clearing denominators gives it in full generality whenever
$M$ is invertible.

For $n\ge1$ take $L=2n+s-1$ and
\[
 I=\{n+s-1,\ldots,2n+s-2\},\qquad J=\{0,\ldots,n-1\}.
\]
Reversing the $n$ columns of $T_L(C)_{I,J}$ produces
$H_n^{[s]}(C)$ and contributes $(-1)^{\binom n2}$. The complementary
minor of $T_L(A)$ has rows $n,\ldots,2n+s-2$ and columns
$0,\ldots,n+s-2$. Reversing its columns produces
$H_{n+s-1}^{[2-s]}(A)$ and contributes
$(-1)^{\binom{n+s-1}{2}}$. Furthermore,
$\sum I+\sum J\equiv ns\pmod2$. Thus the total sign exponent is
\[
 \binom n2+ns+\binom{n+s-1}{2}
 \equiv n+\binom{s-1}{2}\pmod2,
\]
which proves \eqref{eq:recipgeneral} for $n\ge1$.

When $n=0$, the matrix on the right has zeros strictly before its
anti-diagonal and ones on the anti-diagonal. Its determinant is
$(-1)^{\binom{s-1}{2}}$, so the same formula gives one. Finally, apply
\eqref{eq:recipgeneral} with $s=2$, order $n-1$, and the roles of $A,C$
interchanged. This gives \eqref{eq:recipordinary}.
\end{proof}

\section{The Rueppel determinant backbone}\label{sec:backbone}

Define four auxiliary determinant sequences:
\begin{equation}\label{eq:DETB}
 D_n=H_n(r),\qquad E_n=H_n^{[1]}(r),\qquad
 T_n=H_n^{[2]}(r),\qquad B_n=H_n^{[-1]}(r)=H_n(xr(x)).
\end{equation}
All four have value one at $n=0$. The notation $B_n$ is used only for this
auxiliary determinant, not for a power-series coefficient.

\subsection{Unshifted, once-shifted, and twice-shifted determinants}

\begin{theorem}\label{thm:DE}
For $n\ge0$,
\begin{align}
 D_n&=(-1)^{\binom n2},\label{eq:D}\\
 E_n&=(-1)^{\binom{n+1}{2}+\Ones(n)},\label{eq:E}\\
 T_{2m}&=(-1)^m,\qquad T_{2m+1}=0.\label{eq:T}
\end{align}
In addition,
\begin{equation}\label{eq:Erec}
 E_n=D_{\lceil n/2\rceil}E_{\lfloor n/2\rfloor}.
\end{equation}
\end{theorem}
\begin{proof}
The coefficient relations from \eqref{eq:rfunctional} are
\begin{equation}\label{eq:rparity}
 r_{2j+1}=r_j,\qquad r_{2j}=\ind{j=0}.
\end{equation}
Order both rows and columns by even indices first, then odd indices.
Because the same permutation is applied to rows and columns, its two
signs cancel.

For $D_{2m}$, the upper-left block has only its first diagonal entry
nonzero, the lower-right block is zero, and the off-diagonal block is
$(r_{i+j})_{0\le i,j<m}$. Hence
\[
 D_{2m}=(-1)^mD_m^2.
\]
For $D_{2m+1}$, the off-diagonal block has one extra row. Deleting its
first row leaves $(r_{i+j+1})_{0\le i,j<m}$, so
\[
 D_{2m+1}=(-1)^mE_m^2.
\]
Both conclusions use Lemma~\ref{lem:antiblock}.

For $E_n$, the off-diagonal blocks are zero. The even-even block is
an unshifted Rueppel Hankel matrix of order $\lceil n/2\rceil$, and the
odd-odd block is a once-shifted one of order $\lfloor n/2\rfloor$.
This proves \eqref{eq:Erec}. Starting from
$D_0=D_1=E_0=E_1=1$, these formulas simultaneously show that all
$D_n,E_n$ are $\pm1$. They then give
$D_{2m}=D_{2m+1}=(-1)^m$, which is exactly \eqref{eq:D}.
There is no circular use of nonvanishing here: at each order at least
two, the formulas involve smaller orders.

To prove \eqref{eq:E}, it is enough to check that its right side has the
same recurrence and initial value as $E_n$. At $n=2m$, its exponent is
congruent to $m+\Ones(m)$, and the exponent of $D_mE_m$ is
\[
 \binom m2+\binom{m+1}{2}+\Ones(m)
 =m^2+\Ones(m)\equiv m+\Ones(m)\pmod2.
\]
At $n=2m+1$, the exponent of $D_{m+1}E_m$ is congruent to
$\Ones(m)$. Using \eqref{eq:ones-odd}, the proposed expression has
exponent
\[
 (m+1)(2m+1)+\Ones(m)+1-(m\bmod2)
 \equiv\Ones(m)\pmod2.
\]
The value at zero is one, so induction proves the formula.

Finally, the even-even and odd-odd blocks for $T_n$ are both zero.
At even order $2m$, the off-diagonal block is the once-shifted matrix,
with determinant $E_m$. Thus $T_{2m}=(-1)^mE_m^2=(-1)^m$.
At odd order, the two block sizes differ, making the matrix singular.
This proves \eqref{eq:T}.
\end{proof}

A frequently used consequence is
\begin{equation}\label{eq:Tspecial}
 T_{m-1}=D_m\,(m\bmod2)\qquad(m\ge1).
\end{equation}
It follows directly from \eqref{eq:D} and \eqref{eq:T}, including $m=1$.

\subsection{A negative shift detects the number of runs}

\begin{theorem}\label{thm:B}
For every $n\ge1$,
\begin{equation}\label{eq:B}
 B_n=-E_{n-1}\Runs(n-1).
\end{equation}
Consequently, for arbitrary parameters $u,v$,
\begin{equation}\label{eq:linearparam}
 H_n(u+v x r(x))
 =v^{n-1}E_{n-1}\bigl(u-v\Runs(n-1)\bigr),\qquad n\ge1.
\end{equation}
\end{theorem}
\begin{proof}
Put $s_0=0$ and $s_j=r_{j-1}$ for $j\ge1$, so $\sum s_jx^j=xr(x)$.
Then
\[
 s_{2j}=s_j,\qquad s_{2j+1}=\ind{j=0}.
\]
Set $e=\lceil n/2\rceil$, $o=\lfloor n/2\rfloor$. For $n\ge2$, the
parity-ordered matrix for $B_n$ has diagonal blocks given by $B_e$ and
$D_o$, and off-diagonal blocks coupling only their first coordinates.
The first deleted-block determinant is $E_{e-1}$ and the second is
$T_{o-1}$. Lemma~\ref{lem:coupling} therefore gives
\begin{equation}\label{eq:Brec}
 B_n=B_eD_o-E_{e-1}T_{o-1}
 =D_o\bigl(B_e-(o\bmod2)E_{e-1}\bigr).
\end{equation}
The base case is $B_1=0=-E_0\Runs(0)$. Assume the formula holds at
smaller orders. Since $e<n$ for $n\ge2$, \eqref{eq:Brec} becomes
\[
 B_n=-D_oE_{e-1}\bigl(\Runs(e-1)+(o\bmod2)\bigr).
\]
Equation~\eqref{eq:Erec} applied at $n-1$ says
$E_{n-1}=D_oE_{e-1}$. The two digit recurrences
\eqref{eq:runs-even}--\eqref{eq:runs-odd}, separated into even and odd
$n$, give
\[
 \Runs(n-1)=\Runs(e-1)+(o\bmod2).
\]
This completes the induction for \eqref{eq:B}.

For \eqref{eq:linearparam}, changing the constant term of $v x r(x)$
from zero to $u$ changes only the first diagonal entry of its Hankel
matrix. Expansion in that entry gives
\[
 H_n(u+v x r(x))=v^nB_n+u v^{n-1}E_{n-1}.
\]
Insert \eqref{eq:B}. The result is a polynomial identity, so $v=0$
causes no division problem. At $n=1$ the factor $v^0$ is one.
\end{proof}

The specializations $u=1,v=\pm1$ recover the binary-run magnitudes for
$1\pm xr(x)$ that are already part of the earlier Rueppel literature.
The point of the derivation is that the negative shift needed later has
now been evaluated, including its sign and its value at every order.

\section{Proof of the polynomial refinement}\label{sec:mainproof}

Let
\begin{equation}\label{eq:cdef}
 c(x)=\frac1{r(x)}=1-x+x^2-2x^3+3x^4-4x^5+6x^6-\cdots.
\end{equation}
Then $f_t(x)=c(x^2)+tx$. Only the coefficient of $x$ among the
odd-degree coefficients is nonzero. This is what makes the determinant
quadratic in $t$, even when its matrix order is very large.

\begin{proof}[Proof of Theorem~\ref{thm:main}]
The case $n=1$ reads $1=1$, so suppose $n\ge2$ and set
$e=\lceil n/2\rceil$, $o=\lfloor n/2\rfloor$. Reorder by parity.
The even-even block is $(c_{i+j})_{0\le i,j<e}$; the odd-odd block is
$(c_{i+j+1})_{0\le i,j<o}$. The only off-diagonal coupling is $t$ between
their first coordinates. Thus
\begin{align}
 H_n(f_t)
 &=H_e(c)H_o^{[1]}(c)
   -t^2H_{e-1}^{[2]}(c)H_{o-1}^{[3]}(c).\label{eq:ftblock}
\end{align}
Apply Lemma~\ref{lem:reciprocal} four times. The required values are
\begin{align}
 H_e(c)&=(-1)^{e-1}T_{e-1},&
 H_o^{[1]}(c)&=(-1)^oE_o,\label{eq:four1}\\
 H_{e-1}^{[2]}(c)&=(-1)^{e-1}D_e,&
 H_{o-1}^{[3]}(c)&=(-1)^oB_{o+1}.\label{eq:four2}
\end{align}
The last formula also covers $o=1$ by the size-zero part of the
reciprocal-minor lemma. Substitution in \eqref{eq:ftblock} yields
\begin{align}
 H_n(f_t)
 &=(-1)^{n-1}\bigl(T_{e-1}E_o-t^2D_eB_{o+1}\bigr)\notag\\
 &=(-1)^{n-1}D_eE_o\bigl((e\bmod2)+t^2\Runs(o)\bigr)\notag\\
 &=(-1)^{n-1}E_n\bigl(\eps(n)+t^2(\Runs(n)-\eps(n))\bigr).
 \label{eq:ftE}
\end{align}
Here we used \eqref{eq:Tspecial}, \eqref{eq:B}, \eqref{eq:Erec}, and
\eqref{eq:abidentities}, respectively. This already proves the full
polynomial evaluation in a useful alternative form.

It remains to express its sign using the same digit statistics as its
magnitude. Formula~\eqref{eq:E} implies
\[
 (-1)^{n-1}E_n=(-1)^{\binom n2+\Ones(n)+1}.
\]
Write $a=a(n)$, $b=b(n)$, so $\Runs(n)=a+b$ and $b=\eps(n)$.
By Lemma~\ref{lem:digits}, the exponent is congruent to
\[
 1+a+\ceil{(a+b)/2}
 \equiv 1+b+\floor{(a+b)/2}\pmod2.
\]
This is the sign in \eqref{eq:main}. Every step has taken place in the
polynomial ring $\Z[t]$, so no exceptional value of $t$ was excluded.
\end{proof}

\begin{example}[Reading a determinant from its index]
For $n=10$, $g(n)=(1111)_2$, hence $\Runs(n)=4$ and $\eps(n)=1$.
The theorem gives
\[
 H_{10}(f_t)=1+3t^2.
\]
For $n=11$, $g(n)=(1110)_2$, so $\Runs(n)=3$ and $\eps(n)=0$, giving
$H_{11}(f_t)=3t^2$. The index $n=10^{100}$ has $\Runs(n)=122$ and
$\eps(n)=0$, and the result is simply
\[
 H_{10^{100}}(f_t)=122t^2.
\]
The last digit count was computed by the supplied code; no enormous
moment matrix was formed.
\end{example}

\begin{table}[H]
\centering\small
\renewcommand{\arraystretch}{1.1}
\begin{tabular}{@{}r c c r c l@{}}
\toprule
$n$ & Binary $n$ & Gray code $g(n)$ & $\Runs(n)$ & $\eps(n)$ & $H_n(f_t)$\\
\midrule
1 & \texttt{1} & \texttt{1} & 1 & 1 & $1$\\
2 & \texttt{10} & \texttt{11} & 2 & 1 & $-1-t^2$\\
3 & \texttt{11} & \texttt{10} & 1 & 0 & $-t^2$\\
4 & \texttt{100} & \texttt{110} & 2 & 0 & $2t^2$\\
5 & \texttt{101} & \texttt{111} & 3 & 1 & $-1-2t^2$\\
6 & \texttt{110} & \texttt{101} & 2 & 1 & $-1-t^2$\\
7 & \texttt{111} & \texttt{100} & 1 & 0 & $-t^2$\\
8 & \texttt{1000} & \texttt{1100} & 2 & 0 & $2t^2$\\
9 & \texttt{1001} & \texttt{1101} & 3 & 1 & $-1-2t^2$\\
10 & \texttt{1010} & \texttt{1111} & 4 & 1 & $1+3t^2$\\
11 & \texttt{1011} & \texttt{1110} & 3 & 0 & $3t^2$\\
12 & \texttt{1100} & \texttt{1010} & 2 & 0 & $2t^2$\\
13 & \texttt{1101} & \texttt{1011} & 3 & 1 & $-1-2t^2$\\
14 & \texttt{1110} & \texttt{1001} & 2 & 1 & $-1-t^2$\\
15 & \texttt{1111} & \texttt{1000} & 1 & 0 & $-t^2$\\
16 & \texttt{10000} & \texttt{11000} & 2 & 0 & $2t^2$\\
\bottomrule
\end{tabular}
\caption{The first sixteen positive-order determinant polynomials. The sign and both coefficients are read from the index, not from a recurrence in the moments.}
\label{tab:first}
\end{table}


\section{Three further conjectures}\label{sec:three}

\subsection{Reciprocal signs: Conjecture 8}

\begin{theorem}\label{thm:recipsign}
For $n\ge1$ and a parameter $v$,
\begin{equation}\label{eq:qparam}
 H_n\!\left(\frac1{1+v x r(x)}\right)=(-v)^{n-1}E_{n-1}.
\end{equation}
In particular, putting $q_j=H_{j+1}(1/(1+xr(x)))$, we have
\begin{equation}\label{eq:qsign}
 q_j=(-1)^{\binom j2+\Ones(j)}\qquad(j\ge0).
\end{equation}
Therefore
\begin{equation}\label{eq:c8}
 \left|\frac{(-1)^{\binom j2}-q_j}{2}\right|
 =\Ones(j)\bmod2.
\end{equation}
\end{theorem}
\begin{proof}
The coefficients of $1+v x r(x)$ beginning with degree two are
$v r_1,v r_2,\ldots$. Formula~\eqref{eq:recipordinary} therefore gives
\[
 H_n\!\left(\frac1{1+v x r(x)}\right)
 =(-1)^{n-1}v^{n-1}E_{n-1},
\]
proving \eqref{eq:qparam}. At $n=j+1$, \eqref{eq:E} shows that the sign
exponent is
\[
 j+\binom{j+1}{2}+\Ones(j)
 \equiv\binom j2+\Ones(j)\pmod2.
\]
The two signs in the numerator of \eqref{eq:c8} are equal when
$\Ones(j)$ is even and opposite when it is odd.
\end{proof}

The coefficients of $1/(1+xr(x))$ form \oeis{A339422}~\cite{oeisreciprocal}.
The theorem concerns their determinants, not the coefficient sequence
itself. This distinction prevents a common indexing and interpretation
error in numerical searches.

\subsection{A product of signs: Conjecture 23}

\begin{corollary}\label{cor:c23}
For all $j\ge0$,
\begin{equation}\label{eq:c23}
 \frac{1+(-1)^jD_{j+1}E_{j+1}}2
 =\Ones(j+1)\bmod2.
\end{equation}
\end{corollary}
\begin{proof}
With $m=j+1$, \eqref{eq:D} and \eqref{eq:E} give
\[
 D_mE_m=(-1)^{\binom m2+\binom{m+1}{2}+\Ones(m)}
       =(-1)^{m+\Ones(m)}.
\]
Hence $(-1)^jD_{j+1}E_{j+1}=-(-1)^{\Ones(j+1)}$. The expression in
\eqref{eq:c23} is one precisely when $\Ones(j+1)$ is odd.
\end{proof}

Thus Conjectures 8 and 23 are two manifestations of the same once-shifted
Rueppel sign, rather than unrelated coincidences with an OEIS entry.

\subsection{A periodic transform: Conjecture 10}

\begin{theorem}[A two-parameter periodic-block identity]\label{thm:period}
Let $P_{u,v}(x)=u+v x/r(x^2)$. For $m\ge0$,
\begin{equation}\label{eq:Podd}
 H_{2m+1}(P_{u,v})=(-1)^m u v^{2m},
\end{equation}
and, for $m\ge1$,
\begin{equation}\label{eq:Peven}
 H_{2m}(P_{u,v})=
 \begin{cases}
 (-1)^m v^{2m},&m\text{ odd},\\
 0,&m\text{ even}.
 \end{cases}
\end{equation}
The empty determinant is one.
\end{theorem}
\begin{proof}
Again let $c(x)=1/r(x)$. The only even coefficient of $P_{u,v}$ that
may be nonzero is its constant coefficient $u$; its coefficient of
$x^{2j+1}$ is $v c_j$. At order $2m$, parity ordering gives a square
off-diagonal block $v(c_{i+j})$. Lemma~\ref{lem:antiblock} gives
\[
 H_{2m}(P_{u,v})=(-1)^m v^{2m}H_m(c)^2.
\]
By \eqref{eq:recipordinary}, $H_m(c)=(-1)^{m-1}T_{m-1}$, whose square
is one for odd $m$ and zero for even $m$. This proves \eqref{eq:Peven}.

At order $2m+1$, deleting the first row of the rectangular off-diagonal
block leaves $v(c_{i+j+1})$. The rectangular-block identity gives
\[
 H_{2m+1}(P_{u,v})=(-1)^m u v^{2m}H_m^{[1]}(c)^2.
\]
By \eqref{eq:recipgeneral}, $H_m^{[1]}(c)=(-1)^mE_m$, whose square is
one. This also handles $m=0$.
\end{proof}

\begin{corollary}[Conjecture 10]\label{cor:c10}
The Hankel transform, with source index starting at zero, of
$1-x/r(x^2)$ is
\[
 1,-1,-1,0,\ 1,-1,-1,0,\ 1,-1,-1,0,\ldots.
\]
Its generating function is
\begin{equation}\label{eq:periodgf}
 \sum_{j\ge0}H_{j+1}\!\left(1-\frac{x}{r(x^2)}\right)z^j
 =\frac{1-z-z^2}{1-z^4}.
\end{equation}
\end{corollary}
\begin{proof}
Specialize $u=1$, $v=-1$ in Theorem~\ref{thm:period} and separate the
four residue classes of the matrix order. The displayed generating
function is the sum of the resulting period-four blocks.
\end{proof}

\section{Consequences of the polynomial formula}\label{sec:consequences}

\subsection{The exact singular-parameter set}

One reason to retain $t$ instead of specializing immediately is that
\eqref{eq:main} answers a global nonvanishing question at no additional
determinant cost.

\begin{theorem}[All exceptional complex parameters]\label{thm:zeros}
For a fixed $t\in\C$, every leading Hankel determinant of $f_t$ is
nonzero if and only if
\begin{equation}\label{eq:exceptions}
 t\notin\{0\}\ \cup\ 
 \left\{\frac{i}{\sqrt{k}},-\frac{i}{\sqrt{k}}:k=1,2,3,\ldots\right\}.
\end{equation}
Each excluded parameter makes infinitely many positive-order
determinants vanish. In particular, every real $t\ne0$ gives nonzero
determinants of every order.
\end{theorem}
\begin{proof}
Use $a=a(n)$ and $b=b(n)$ from \eqref{eq:ab}. Up to a sign, the
polynomial is $b+a t^2$. For $n\ge1$, $(a,b)\ne(0,0)$.
If $b=0$, then $a\ge1$, and the only root is $t=0$.
If $b=1$ and $a=0$, the polynomial is one. If $b=1$ and $a=k\ge1$,
its roots are precisely $\pm i/\sqrt{k}$.

It remains to show that none of these possibilities is merely formal.
Every nonnegative integer occurs as a Gray code: given its digits,
recover the binary digits by cumulative xor from the highest position.
For each $k\ge1$, infinitely many Gray codes have lowest bit one and
exactly $k$ ones above that bit. They yield infinitely many indices
with the factor $1+k t^2$. Similarly, infinitely many nonzero Gray
codes have lowest bit zero, yielding roots at $t=0$. The converse
follows directly from the exhaustive cases for $b+a t^2$.
\end{proof}

Nonvanishing here is not positivity. For example,
$H_2(f_t)=-1-t^2<0$ for real $t$, so these moments do not define a
positive-definite real Hankel form. The result is an exact
nondegeneracy statement.

\subsection{Polynomial and signed-value distributions}

\begin{theorem}[Dyadic polynomial distribution]\label{thm:distribution}
Fix $N\ge1$. Among indices $1\le n<2^N$, the pair $(a(n),b(n))=(k,b)$,
where $0\le k\le N-1$ and $b\in\{0,1\}$, occurs
\begin{equation}\label{eq:distribution}
 \binom{N-1}{k}-\ind{k=0,\ b=0}
\end{equation}
times. All indices with that pair have the same determinant polynomial,
namely
\begin{equation}\label{eq:polyclasses}
 H_n(f_t)=
 \begin{cases}
 (-1)^{1+\lfloor k/2\rfloor}k t^2,&b=0,\\
 (-1)^{\lfloor(k+1)/2\rfloor}(1+k t^2),&b=1.
 \end{cases}
\end{equation}
\end{theorem}
\begin{proof}
The Gray-code map permutes the $N$-bit strings. With the lowest bit
fixed, there are $\binom{N-1}{k}$ ways to choose the $k$ higher ones.
The all-zero string corresponds to the excluded index $n=0$.
The two signs in \eqref{eq:polyclasses} are the main theorem's sign,
written separately for $b=0$ and $b=1$.
\end{proof}

A compact form of this distribution is the bivariate identity
\begin{equation}\label{eq:distributiongf}
 \sum_{1\le n<2^N}u^{a(n)}v^{b(n)}
 =(1+v)(1+u)^{N-1}-1.
\end{equation}
It describes the constant and quadratic coefficients before the
explicit sign is attached; it is not a generating function of the
original power-series coefficients.

\begin{corollary}[Exact signed magnitudes at $t=1$]\label{cor:signeddist}
For $1\le k\le N$, put $s_k=(-1)^{\lfloor k/2\rfloor}$. Among
$1\le n<2^N$,
\begin{align}
 \#\{n:H_n(f_1)=s_k k\}&=\binom{N-1}{k-1},\label{eq:positivecount}\\
 \#\{n:H_n(f_1)=-s_k k\}&=\binom{N-1}{k}.\label{eq:negativecount}
\end{align}
In particular, $|H_n(f_1)|=k$ occurs $\binom Nk$ times.
\end{corollary}
\begin{proof}
For magnitude $k$, the case $b=1$ has $a=k-1$ and sign $s_k$;
the case $b=0$ has $a=k$ and the opposite sign. Add the two counts and
use Pascal's identity. Binomial coefficients outside their usual
range are understood to be zero.
\end{proof}

The largest magnitude on this interval is $N$, attained once, at the
index with all $N$ Gray-code digits equal to one. Its binary digits
alternate, and its value is
\begin{equation}\label{eq:maxindex}
 n=\floor{2^{N+1}/3},\qquad
 H_n(f_1)=(-1)^{\lfloor N/2\rfloor}N.
\end{equation}
Thus the determinants grow at most logarithmically in their matrix
order, despite the much larger coefficients from which they are built.

\subsection{Dyadic sums with the parameter retained}

\begin{theorem}\label{thm:sums}
For $N\ge2$ and arbitrary $t$,
\begin{equation}\label{eq:signedsum}
 \sum_{1\le n<2^N}H_n(f_t)
 =\operatorname{Re}(1+i)^N
 -2(N-1)t^2\operatorname{Re}(1+i)^{N-2}.
\end{equation}
At $N=1$ the sum is one. For real $t$ and all $N\ge1$,
\begin{equation}\label{eq:abssum}
 \sum_{1\le n<2^N}|H_n(f_t)|
 =2^{N-1}\bigl(1+(N-1)t^2\bigr).
\end{equation}
The real parts in \eqref{eq:signedsum} are integers; the identity is
an identity of polynomials in $t$.
\end{theorem}
\begin{proof}
Set $M=N-1$. Let $\sigma_0(k)=(-1)^{1+\lfloor k/2\rfloor}$ and
$\sigma_1(k)=(-1)^{\lfloor(k+1)/2\rfloor}$. Their four-periodic patterns
show that
\[
 \sigma_1(k)=\operatorname{Re}((1+i)i^k),\qquad
 \sigma_0(k)+\sigma_1(k)=-2\operatorname{Im}(i^k).
\]
The missing pair $(0,0)$ contributes zero, so
Theorem~\ref{thm:distribution} gives a constant coefficient
\[
 \sum_{k=0}^{M}\binom Mk\sigma_1(k)
 =\operatorname{Re}\bigl((1+i)(1+i)^M\bigr)
 =\operatorname{Re}(1+i)^N.
\]
The coefficient of $t^2$ is
\begin{align*}
 -2\operatorname{Im}\sum_{k=0}^{M}k\binom Mk i^k
 &=-2\operatorname{Im}\bigl(Mi(1+i)^{M-1}\bigr)\\
 &=-2M\operatorname{Re}(1+i)^{M-1}.
\end{align*}
This proves \eqref{eq:signedsum} when $M\ge1$; the case $N=1$ is direct.

For real $t$, the unsigned polynomial is $b+a t^2\ge0$.
Summing over all pairs gives
\[
 \sum_{k=0}^{M}\binom Mk\bigl(k t^2+(1+k t^2)\bigr)
 =2^M+M2^M t^2,
\]
using $\sum k\binom Mk=M2^{M-1}$ for $M\ge1$. The case $M=0$ is
again immediate. This is \eqref{eq:abssum}.
\end{proof}

\section{A printed counterexample and a signed replacement}\label{sec:correction}

The four affirmative results above are the principal conclusions.
There is also a small but important correction in the same source.
We treat it separately because finding an omitted sign or absolute
value is not comparable in significance to a structural theorem.

\subsection{What fails in Conjecture 22}

Let
\begin{equation}\label{eq:complementdefs}
 C_m=\det(1-r_{i+j})_{0\le i,j<m},\qquad
 J_m=\det(r_{i+j+1}-r_{i+j})_{0\le i,j<m}.
\end{equation}
The formula printed as Conjecture 22 would imply, in matrix-size
notation,
\begin{equation}\label{eq:printed22}
 J_m^2=|C_{m+1}|-|C_m|\qquad(m\ge1).
\end{equation}
At $m=3$, however,
\[
 C_3=\det\begin{pmatrix}0&0&1\\0&1&0\\1&0&1\end{pmatrix}=-1,
 \qquad
 C_4=\det\begin{pmatrix}
 0&0&1&0\\0&1&0&1\\1&0&1&1\\0&1&1&1
 \end{pmatrix}=0.
\]
The last row of the second matrix is the sum of its first two rows.
On the other hand,
\[
 J_3=\det\begin{pmatrix}0&-1&1\\-1&1&-1\\1&-1&0\end{pmatrix}=1.
\]
Thus \eqref{eq:printed22} asserts $1=-1$. In the source's indexing,
this is the counterexample $j=2$. Direct checks at $m=1,2$ satisfy
the formula, so this is its first failure. The assertion is therefore
false as printed; an additional absolute value or a different sign
convention would constitute a different assertion.

\subsection{A general moment identity}

There is a natural exact replacement that avoids guessing which signs
to suppress. In fact it applies to arbitrary moment sequences.

\begin{theorem}[Complement--difference square identity]\label{thm:square}
Let $(a_j)_{j\ge0}$ be any sequence, and define
\[
 d_m=\det(a_{i+j}),\qquad c_m=\det(1-a_{i+j}),\qquad
 j_m=\det(a_{i+j+1}-a_{i+j}),
\]
where all matrices have order $m$. Then for $m\ge1$,
\begin{equation}\label{eq:generalsquare}
 j_m^2=(-1)^m\bigl(d_m c_{m+1}+d_{m+1}c_m\bigr).
\end{equation}
This polynomial identity requires no nonvanishing assumptions on the
determinants. In the Rueppel case it becomes
\begin{equation}\label{eq:rueppelsquare}
\boxed{\quad J_m^2=D_mC_m+D_{m+1}C_{m+1}.\quad}
\end{equation}
\end{theorem}
\begin{proof}
We first work with generic moments over a rational-function field,
where the needed inverses exist. Put
\[
 A_m=(a_{i+j})_{0\le i,j<m},\quad
 v=(a_m,\ldots,a_{2m-1})^T,\quad
 K_m=\one^TA_m^{-1}\one.
\]
The matrix determinant lemma gives
\begin{equation}\label{eq:Cnorm}
 c_m=(-1)^m d_m(1-K_m).
\end{equation}
The next Hankel matrix is
\[
 A_{m+1}=\begin{pmatrix}A_m&v\\v^T&a_{2m}\end{pmatrix},
\]
whose Schur complement is
$s=a_{2m}-v^TA_m^{-1}v=d_{m+1}/d_m$. Its block inverse gives
\begin{equation}\label{eq:Kdiff}
 K_{m+1}-K_m
 =\frac{(1-v^TA_m^{-1}\one)^2}{s}.
\end{equation}
To identify the numerator, consider the bordered matrix
\[
 M=\begin{pmatrix}
 a_0&a_1&\cdots&a_m\\
 a_1&a_2&\cdots&a_{m+1}\\
 \vdots&\vdots&&\vdots\\
 a_{m-1}&a_m&\cdots&a_{2m-1}\\
 1&1&\cdots&1
 \end{pmatrix}.
\]
Its determinant by a Schur complement is
$d_m(1-v^TA_m^{-1}\one)$. Alternatively, for columns $m,m-1,\ldots,1$,
subtract the preceding column. The last row becomes $(1,0,\ldots,0)$;
expanding there leaves the first-difference Hankel determinant with
sign $(-1)^m$. Therefore
\[
 d_m(1-v^TA_m^{-1}\one)=(-1)^m j_m.
\]
Substitution into \eqref{eq:Kdiff} yields
\[
 K_{m+1}-K_m=\frac{j_m^2}{d_md_{m+1}}.
\]
Now use \eqref{eq:Cnorm} at the two consecutive orders and rearrange:
\[
 j_m^2=(-1)^m(d_mc_{m+1}+d_{m+1}c_m).
\]
Both sides are polynomials with integer coefficients in finitely many
moments. Equality for generic moments proves the polynomial identity,
so specialization is valid even when some determinants vanish and over
any commutative ring.

For the Rueppel sequence, $D_{m+1}=(-1)^mD_m$. Inserting this relation
into \eqref{eq:generalsquare} gives \eqref{eq:rueppelsquare}.
\end{proof}

At the counterexample above, the corrected formula reads
$J_3^2=D_3C_3+D_4C_4=(-1)(-1)+(1)(0)=1$, as required.
We do not claim that simply adding an outer absolute value to the
printed formula has been proved here; the precise replacement established
is the signed polynomial identity \eqref{eq:rueppelsquare}.

\section{Reproducibility, status, and remaining questions}\label{sec:status}

\subsection{What was checked computationally}

The mathematical proofs above apply at every order. Computation serves
as an independent audit of the indexing, signs, polynomial specializations,
and implementation; it is not a substitute for those proofs.

The archive contains a standard-library Python module, a separate test
suite, an exact verification script, and the outputs of an actual run.
The determinant engine is fraction-free Bareiss elimination with row
pivoting. Every division is checked for exactness. Thus neither a small
floating-point residual nor an approximate zero is accepted as evidence.
The engine does not call any of the binary formulas when evaluating a
matrix determinant.

\begin{center}
\small
\begin{tabular}{@{}p{.55\textwidth}p{.35\textwidth}@{}}
\toprule
Audit & Executed scope\\
\midrule
Rueppel-related determinants & 1,192 exact determinant evaluations.\\
Consecutive orders & All orders $0$ through $64$ for each tested family.\\
Selected larger orders & $65$, $127$, $128$ for the seven indicated base families.\\
Parameter values for $f_t$ & $t=-2,0,1,2$.\\
General square identity & 500 checks on randomly generated integer moments.\\
Digital recurrence for $E_n$ & Every $0\le n<2^{16}$, plus 1,000 random 256-bit indices.\\
Dyadic distributions and sums & Distribution at $N=16$; sums for $1\le N\le16$, $t=0,1,2$.\\
Independent determinant audit & Leibniz permutation sums on 140 random matrices, sizes $0$ through $6$.\\
Reciprocal-minor identity audit & Random integer series and shifts $1$ through $5$, including negative complementary shifts.\\
\bottomrule
\end{tabular}
\end{center}

The seven families receiving the larger-order checks are $D,E,T,B$,
$f_1$, $1-x/r(x^2)$, and $1/(1+xr(x))$. Complement determinants were
also computed at order $65$ to check the square identity through order
$64$. The complete verification run used Python 3.13.5 and took about
7.29 seconds in the execution environment; this is a record of that run,
not a portable timing guarantee. All checks passed.

Two implementation details are particularly important. First, the
size-zero determinant is handled explicitly and is never confused
with the first term of a Hankel transform. Second, singular matrices
and zero leading pivots are tested intentionally. The periodic family
contains infinitely many zero determinants, so a routine that silently
assumes nonzero pivots would be unsuitable.

From the archive directory, the two reproduction commands are
\begin{lstlisting}[language={}]
python -m unittest discover -s tests -v
python code/verify.py
\end{lstlisting}
The default verification output overwrites files under \texttt{data/}.
Larger finite audits can be requested, for example,
\begin{lstlisting}[language={}]
python code/verify.py --max-size 96 --large-size 192 --digital-bits 18
\end{lstlisting}
This last command is an optional example, not a run claimed in the
reported results. No external Python packages, internet access, Wolfram
kernel, or proof assistant are required by the supplied code.

\subsection{Source and novelty audit}

The primary problem source is \arxiv{2107.00442v2}, dated 4 July 2021.
The statements and formulas were checked against the PDF, specifically
printed pages 6--8 and 14, as well as the HTML representation.
The arXiv abstract page's current version history was checked on
19 September 2026. The earlier paper by Allouche, Han, and
Shallit was consulted through \arxiv{2006.08909v2}; its journal version
is listed in the bibliography. The relevant OEIS entries were also
checked on that date.

The present proofs establish the stated identities mathematically,
subject to ordinary independent scrutiny of a new manuscript. They do
not establish bibliographic priority. In particular, the evaluation of
$D_n$, recurrences for $E_n$, and binary-run phenomena for $1\pm xr(x)$
are prior material. The parameter formula, the explicit deductions of
the four targeted 2021 assertions, and their collected consequences are
the outputs of this investigation. A focused search did not identify a
later paper resolving those exact assertions, but unindexed proofs,
unpublished communications, or equivalent formulations may exist.
No claim is made that an OEIS record or the author's paper has been
updated to reflect this manuscript.

The correction to Conjecture 22 is deliberately described as a
counterexample to the \emph{printed statement}. Its small size makes a
typographical or sign-convention explanation plausible; nothing here
purports to infer the author's intended corrected formula.

\subsection{Boundaries of the result}

The proofs do not resolve every conjecture in the 2021 paper.
In particular, no conclusion is asserted here about its other
parameterized families or its period-eight assertion modulo two.

The method suggests two concrete extensions. One is to evaluate higher
shifts $H_n^{[s]}(r)$ uniformly in both $n$ and $s$, which would make the
reciprocal-minor lemma more broadly effective. Another is to replace
the perturbation $tx$ by a polynomial with several odd-degree terms.
Parity ordering would then produce a finite-rank coupling, but its
minors would require more than the four determinant sequences used
here. These are directions for further work, not results asserted by
this manuscript.

\appendix
\section{Direct evaluation and complexity}\label{sec:algorithm}

The main formula is not merely a recurrence. It is an explicit function
of the binary digits of the matrix order. The following code returns
the constant and quadratic coefficients of the determinant polynomial.
It is also included, with input checks and the other family formulas,
in \texttt{code/rueppel.py}.

\begin{lstlisting}
def parameter_coefficients(n: int) -> tuple[int, int]:
    if not isinstance(n, int) or isinstance(n, bool) or n < 0:
        raise ValueError("The index must be a nonnegative integer.")
    if n == 0:
        return 1, 0
    gray = n ^ (n >> 1)
    a = (gray >> 1).bit_count()
    b = gray & 1
    exponent = 1 + b + (a + b) // 2
    sigma = -1 if exponent % 2 else 1
    return sigma * b, sigma * a
\end{lstlisting}

If the input index has $L$ binary digits, xor, shifting, and population
count need $O(L)$ bit work in a straightforward implementation.
The output coefficients have $O(\log L)$ bits. Thus the determinant
polynomial can be evaluated using work linear in the input bit length,
without generating moments or building an $n\times n$ matrix.
The displayed Python implementation stores intermediate integers of
$O(L)$ bits; it is not a claim of constant-space computation for
arbitrary-precision inputs.

For comparison, the independent matrix audit uses a cubic number of
integer arithmetic operations per determinant, with bit cost depending
on the intermediate entries. Its purpose is verification from the
definition. Replacing it by the binary formula inside the audit would
make the main comparison circular.

\section{Archive contents}

\begin{center}
\small
\begin{tabular}{@{}p{.44\textwidth}p{.48\textwidth}@{}}
\toprule
File & Purpose\\
\midrule
\texttt{rueppel\_hankel.tex}, \texttt{tables.tex} & Complete article source and generated table.\\
\texttt{rueppel\_hankel.pdf} & Compiled article.\\
\texttt{code/rueppel.py} & Exact formulas, formal reciprocals, and determinant engine.\\
\texttt{code/verify.py} & Reproduction script for the recorded experiments.\\
\texttt{code/make\_table.py} & Regenerates the first-values table.\\
\texttt{tests/test\_rueppel.py} & Independent small-size and identity tests.\\
\texttt{data/exact\_determinants.csv} & Direct matrix evaluations through order 64.\\
\texttt{data/polynomial\_coefficients.csv} & Binary data and determinant polynomials through order 1024.\\
\texttt{data/dyadic\_distribution.csv} & The coefficient distribution at $N=16$.\\
\texttt{data/dyadic\_sums.json} & Exact sums for the tested dyadic ranges.\\
\texttt{data/verification\_summary.json} & Machine-readable audit summary and counterexample.\\
\texttt{data/verification\_log.txt} & Output of the executed verification run.\\
\texttt{README.md}, \texttt{Makefile} & Reading and build instructions.\\
\texttt{SOURCE\_AUDIT.md} & Source locations and priority caveats.\\
\bottomrule
\end{tabular}
\end{center}

\begin{thebibliography}{9}
\bibitem{barry2021}
Paul Barry,
\emph{Conjectures and results on some generalized Rueppel sequences},
\arxiv{2107.00442v2}, 4 July 2021.
\href{https://doi.org/10.48550/arXiv.2107.00442}{doi:10.48550/arXiv.2107.00442}.

\bibitem{ahs}
Jean-Paul Allouche, Guo-Niu Han, and Jeffrey Shallit,
\emph{On some conjectures of P. Barry},
Journal of Number Theory \textbf{228} (2021), 108--132.
\href{https://doi.org/10.1016/j.jnt.2021.03.015}{doi:10.1016/j.jnt.2021.03.015}.
Preprint: \arxiv{2006.08909v2}, 24 June 2020.

\bibitem{barry2020}
Paul Barry,
\emph{Some observations on the Rueppel sequence and associated Hankel determinants},
\arxiv{2005.04066}, 2020.

\bibitem{oeisrueppel}
The OEIS Foundation Inc.,
\oeis{A036987}, the Rueppel sequence. Accessed 19 September 2026.

\bibitem{oeisruns}
The OEIS Foundation Inc.,
\oeis{A005811}, number of runs in the binary expansion of $n$.
Accessed 19 September 2026.

\bibitem{oeisparity}
The OEIS Foundation Inc.,
\oeis{A268411}, parity of the number of runs of ones in the binary
representation of $n$. Accessed 19 September 2026.

\bibitem{oeisreciprocal}
The OEIS Foundation Inc.,
\oeis{A339422}, coefficients of $1/(1+\sum_{j\ge0}x^{2^j})$.
Accessed 19 September 2026.
\end{thebibliography}
\end{document}
